\documentclass[acmlarge,screen,nonacm]{acmart}

% --- packages not already loaded by acmart ---
\usepackage{tabularx}

% suppress the ACM reference block (also implied by nonacm)
\settopmatter{printacmref=false}
\setlength{\footskip}{12pt}

% --- float placement tuning -------------------------------------------
% The defaults are conservative (topfraction 0.7, totalnumber 3), which
% causes tall figures restricted to [t] to be deferred and pile up at the
% end of the document. Loosening these lets figures sit near their text.
\renewcommand{\topfraction}{0.9}        % up to 90% of a page top can be floats
\renewcommand{\bottomfraction}{0.9}     % same for the bottom
\renewcommand{\textfraction}{0.07}      % require very little text per page
\renewcommand{\floatpagefraction}{0.7}  % a float page must be >=70% full
\setcounter{topnumber}{3}               % allow more floats at the top
\setcounter{bottomnumber}{3}            % and at the bottom
\setcounter{totalnumber}{5}             % and more floats overall per page


%=====================================================================
\begin{document}

\title{Tactile Hands: A Force-Sensing Robotic Hand for Safe Grasping}

\author{Bryce Hackel}
\authornote{Project lead.}
\affiliation{%
  \institution{University of California, San Diego}
  \city{La Jolla}\state{California}\country{USA}}
\email{bhackel@ucsd.edu}


\author{Ali El Lahib}
\affiliation{%
  \institution{University of California, San Diego}
  \city{La Jolla}\state{California}\country{USA}}
\email{aellahib@ucsd.edu}



\author{Thomas Nghiem}
\affiliation{%
  \institution{University of California, San Diego}
  \city{La Jolla}\state{California}\country{USA}}
\email{tnghiem@ucsd.edu}

\author{Shree Venkatesh}
\affiliation{%
  \institution{University of California, San Diego}
  \city{La Jolla}\state{California}\country{USA}}
\email{s1venkatesh@ucsd.edu}

\author{Sidath Daham Wijesinghe}
\affiliation{%
  \institution{University of California, San Diego}
  \city{La Jolla}\state{California}\country{USA}}
\email{swijesinghe@ucsd.edu}

\renewcommand{\shortauthors}{Triton Droids: Tactile Hands}

\begin{abstract}
Robotic hands that grasp without sensing how hard they squeeze can crush
fragile objects, yet no open-source hand under \$200 offers fingertip force
sensing. We present Tactile Hands, a sub-\$200, four-fingered robotic hand that
adds calibrated fingertip force-sensing resistors (FSRs) to the open-source
AmazingHand. A custom ESP32-S3 board reads four FSRs through voltage dividers
and streams calibrated forces to a ROS~2 host, where a force-aware controller
halts finger closure at a safe threshold. Calibration reaches better-than-1~N
accuracy. Mounted on a teleoperated arm and driven in VR, the same controller
grasps a thin-walled cup without crushing it, where a force-blind grasp collapses
it. This is a low-cost foundation for tactile manipulation.
\end{abstract}

\keywords{tactile sensing, force-sensing resistors, robotic hand, grasping,
calibration, teleoperation, ROS~2, low-cost robotics}

\maketitle

%=====================================================================
\section{Introduction}
%=====================================================================
A robot arm can be precise to the millimeter and still be dangerous, because
position tells it nothing about force. A hand that cannot feel how hard it is
gripping will crush a paper cup as readily as it lifts a steel block.
Humans avoid this through a dense array of tactile receptors in the fingertips
that report contact force in real time and let us modulate our grip
unconsciously. Bringing the same capability to robots, fingertip force
sensing tightly coupled to the grip controller, is what separates a hand that
``closes until it stops'' from a hand that ``closes until it feels enough.''

High-end research hands such as the Shadow Dexterous Hand or sensorized hands
built around GelSight~\cite{yuan2017gelsight} and DIGIT~\cite{lambeta2020digit}
tactile sensors demonstrate the value of touch, but they cost thousands of
dollars and are out of reach for most students and small labs.
Low-cost open-source hands have begun to close this gap on the mechanical side:
the LEAP Hand~\cite{shaw2023leap} and the AmazingHand~\cite{pollen2025amazinghand}
are anthropomorphic, 3D-printable, and affordable. However, these affordable
platforms ship without fingertip force sensing. To our knowledge, there is no
open-source robotic hand available for under roughly \$200 that provides
calibrated fingertip force measurement and uses it to grasp safely. This is the
gap our project targets.

Why is this capability missing at low cost? The sensors that make grip force
legible to a robot tend to be expensive. Vision-based tactile sensors such as
GelSight~\cite{yuan2017gelsight} and DIGIT~\cite{lambeta2020digit} embed a camera
and illumination behind a deformable gel and recover rich contact geometry, but
each fingertip then needs its own camera and per-frame computation, workable for
a single research fingertip but prohibitive across four fingers on a student budget.
Force-sensing resistors (FSRs) sit at the opposite extreme: a few dollars each, a
single analog channel per sensor, and no per-finger compute. Their drawback is
signal quality: FSRs are nonlinear, hysteretic, and vary unit to unit, so a raw
reading is not a force. We treated this as an engineering problem rather than a
fundamental limit. With a controlled calibration procedure and a simple
per-sensor model, inexpensive FSRs can be made accurate enough (better than
1~N) to close a useful safety loop, which is all that grasping a fragile object
without crushing it actually requires.

We present \textbf{Tactile Hands}, a four-fingered robotic hand built on the
open-source AmazingHand mechanical platform and augmented with a fingertip
force-sensing subsystem of our own design. Each fingertip carries a
force-sensing resistor (FSR); a custom ESP32-S3 board reads all four sensors
through voltage dividers, converts raw ADC readings to Newtons using a
calibrated power-law model, and streams the values over USB to a ROS~2 host. A
force-aware grasp controller continuously compares each fingertip force against
a safe threshold and stops closing a finger the moment that threshold is
crossed, allowing the hand to wrap a thin-walled paper cup and hold it without
collapsing it. The complete bill of materials remains under \$200. That figure
covers the hand itself driven by a laptop webcam; the Valve Index controllers
and ARCTOS arm are stretch additions on top.

Two design choices matter more than the parts list.
First, all sensing and actuation are mediated by a single ROS~2 interface:
calibrated forces are published on one topic and servo targets are read from
another, so the grasp controller never talks to hardware directly and any input
source (a fixed grasp trajectory, a webcam, or a VR controller) can drive the
hand through the same contract. This let five subsystems (electronics/firmware,
calibration, simulation, vision, and VR) be built and tested in parallel against
a stable interface rather than serialized one behind another. Second, the
force-aware controller is deliberately minimal: a per-finger threshold that
freezes a finger the instant its measured force crosses a safe limit. We show
that this minimal loop, running on calibrated low-cost sensors, is enough to make
the hand stop squeezing at a safe force, and that the same guarantee keeps
holding once the hand is mounted on a teleoperated arm and driven by a human.

Beyond the core force-sensing minimum viable product (MVP), we built three
input and simulation extensions that share the same ROS~2 backbone: a
webcam-based hand-pose pipeline that mirrors a human operator's finger curl onto
the hand, a physics simulation of the AmazingHand for kinematics validation and
future sim-to-real work, and a virtual-reality (VR) teleoperation pipeline that
streams 6-DoF controller pose into ROS~2. Most recently, we mounted the hand on
an open-source ARCTOS robotic arm~\cite{arctos2024} and teleoperated its
end-effector from VR controllers. The force-aware controller keeps
running while a human drives the arm: under VR teleoperation the hand grasps a
thin-walled cup and halts at the safe limit, whereas the identical
teleoperated motion with force feedback disabled buckles the cup wall. We
further added vision-guided, multi-object grip selection: a fine-tuned YOLO26
object detector recognizes the target object in a webcam image and selects an
appropriate per-object grip force automatically, which the same force-aware loop
then enforces.

\paragraph{Contributions.} This paper makes the following contributions:
\begin{itemize}
  \item A \textbf{sub-\$200, open hardware design} that adds fingertip force
        sensing to a low-cost, 3D-printable robotic hand, including a custom
        ESP32-S3 sensing board with per-channel voltage dividers.
  \item A \textbf{repeatable FSR calibration procedure} using stacked 3D-printed
        reference weights and a power-law fit that maps raw ADC counts to force
        with better-than-1~N accuracy (max error 25.7~g $\approx$ 0.25~N).
  \item A \textbf{force-aware grasp controller} on a ROS~2 host that fuses the
        calibrated fingertip forces with finger position commands to grasp and
        release a fragile object without crushing it.
  \item A \textbf{demonstration that the force-aware safety guarantee survives
        teleoperation}: under VR control of the arm-mounted hand, force-thresholded
        grasping holds a fragile cup that an otherwise-identical force-blind
        grasp visibly collapses---isolating the contribution of the tactile
        channel.
  \item A set of \textbf{teleoperation and simulation extensions}
        (webcam hand-pose mirroring, physics simulation, and VR controller
        streaming) that share one ROS~2 interface, culminating in a
        VR-teleoperated, arm-mounted hand and a vision-guided, multi-object
        grip-selection pipeline driven by a fine-tuned YOLO26 detector.
\end{itemize}

The remainder of the paper reviews related work
(Section~\ref{sec:related}), details the system design and each subsystem
(Section~\ref{sec:technical}), reports our milestone history including
mid-quarter revisions and the problems we encountered
(Section~\ref{sec:milestones}), and concludes with future work
(Section~\ref{sec:conclusion}).

%=====================================================================
\section{Related Works}
\label{sec:related}
%=====================================================================
Our project sits at the intersection of low-cost robot hands, tactile sensing,
and teleoperation. We summarize the most relevant prior work in each area and
explain how Tactile Hands relates to it.

\paragraph{Low-cost and open-source robotic hands.}
The LEAP Hand~\cite{shaw2023leap} demonstrated that an anthropomorphic,
dexterous hand can be built for a few hundred dollars from off-the-shelf servos
and 3D-printed parts, and it has become a popular platform for robot learning.
The AmazingHand~\cite{pollen2025amazinghand} pushes cost down further with a
fully 3D-printed, tendon- and linkage-driven design powered by small serial-bus
servos; we adopt it as our mechanical base. These platforms prioritize low cost
and reproducibility, but they ship without fingertip force sensing. Our
work keeps their affordability while adding the missing tactile channel.

\paragraph{Tactile sensing technologies.}
Tactile sensing for robots spans a wide design space, surveyed by Dahiya et
al.~\cite{dahiya2010tactile}. High-resolution vision-based sensors such as
GelSight~\cite{yuan2017gelsight} and DIGIT~\cite{lambeta2020digit} recover both
contact geometry and force at the cost of an embedded camera, illumination, and
significant compute per fingertip: excellent fidelity, but poorly matched to a
sub-\$200 budget across four fingers. At the opposite end, force-sensing
resistors (FSRs) are inexpensive piezoresistive sensors whose resistance drops
nonlinearly with applied normal force~\cite{interlink2024fsr}. FSRs are noisy,
hysteretic, and require per-sensor calibration, but they cost on the order of a
few dollars and need only a single ADC channel each. We choose FSRs precisely
because they make four-fingertip force sensing economically feasible, and we
address their nonlinearity with an explicit calibration model.

\paragraph{Force-aware and safe grasping.}
A long line of work shows that closing the loop on contact force improves
grasping robustness and prevents damage to fragile objects. Learning-based
systems such as OpenAI's in-hand manipulation
work~\cite{andrychowicz2020dexterity} achieve remarkable dexterity but rely on
extensive simulation and high-end hardware. Our controller is deliberately
simple, a thresholded force limit on each finger, and it shows that even a
minimal force feedback loop on cheap sensors yields a meaningful safety
guarantee.

\paragraph{Teleoperation.}
Vision- and VR-based teleoperation map a human operator's motion onto a robot.
DexPilot~\cite{handa2020dexpilot} and AnyTeleop~\cite{qin2023anyteleop}
teleoperate dexterous hand-arm systems from camera observations of the
operator's hand. For lightweight pose capture we use Google's MediaPipe
Hands~\cite{zhang2020mediapipe}, which provides real-time 21-landmark hand
tracking from a single webcam. For VR we stream controller pose using Valve's
OpenVR~\cite{valve2024openvr} interface, building on the
\texttt{unity\_ros\_teleoperation} framework~\cite{leggedrobotics_teleop}.

\paragraph{Object recognition for grip selection.}
Choosing how firmly to grasp an object first requires recognizing what it is.
The YOLO (``You Only Look Once'') family of single-stage detectors is the
standard tool for fast, real-time object detection, and the latest release,
Ultralytics YOLO26~\cite{ultralytics_yolo26}, is engineered for efficient,
low-latency inference on edge hardware. Because the appropriate grip force
depends on the object---a full bottle must be held more firmly than an empty
one---we fine-tune YOLO26 on a small custom-annotated dataset of our test
objects and use the predicted class to look up a per-object force threshold
(Section~\ref{sec:multiobject}), turning recognition into an automatic
grip-force selector that feeds the existing force-aware loop.

\paragraph{Simulation and sim-to-real.}
Physics simulators such as MuJoCo~\cite{todorov2012mujoco} and NVIDIA Isaac
Gym~\cite{makoviychuk2021isaacgym} are standard tools for validating robot
kinematics and training control policies before deployment. Transferring such
policies to hardware remains challenging due to the reality
gap~\cite{zhao2020sim2real}.
We bring up the AmazingHand in simulation to verify finger
kinematics and to lay the groundwork for sim-to-real in-hand manipulation,
while keeping all real-time control on ROS~2~\cite{macenski2022ros2} so that
simulation, vision, and VR subsystems can be developed in parallel against the
same interface.

%=====================================================================
\section{Technical Material}
\label{sec:technical}
%=====================================================================

\subsection{System Architecture}
The system is organized as a sensing-and-actuation hand connected to a ROS~2
host. On the hand, four FSRs are read by an ESP32-S3 microcontroller through
voltage dividers; the firmware samples, converts to Newtons, and streams the
forces over USB serial. On the host, a serial-bridge node publishes the forces
on a ROS~2 topic and forwards servo commands back to the hand. The same host
runs the input subsystems---vision-based hand-pose estimation, VR controller
streaming---and the simulation, all of which produce target finger (and, for the
arm, end-effector) commands. The force-aware grasp controller is the arbiter
between desired motion and measured force: it allows commanded motion only while
fingertip forces remain below a safe limit. Figure~\ref{fig:block} shows this
control logic.

\begin{figure}[htbp]
\centering
\includegraphics[width=0.9\linewidth]{figures/control_loop.png}
\caption{Force-aware grasp control loop. Target finger positions (from a fixed
grasp trajectory or a teleoperation input) and live FSR forces are compared at
every cycle. If a finger's force exceeds the safe threshold, that finger's
motion is restricted and a red status LED lights; otherwise the servo command
is issued.}
\label{fig:block}
\end{figure}

\subsection{Mechanical Platform}
The hand is based on the open-source AmazingHand~\cite{pollen2025amazinghand}, a
four-fingered, fully 3D-printed design. Each finger is actuated by two Feetech
SCS0009 serial-bus micro servos (eight in total), giving four independently
controllable fingers.
We retained the original four-finger morphology rather than expanding to five:
four fingers are sufficient to demonstrate force-aware power and pinch grasps on
our test objects, simplify the control problem, and keep the hand compact enough
to mount on a small arm. The fingertips were modified to host an FSR pad and to
route its leads back toward the palm (Figure~\ref{fig:mvphand}). The entire hand,
prototype board, and status LEDs are enclosed in a 3D-printed shell printed on
a Bambu Lab A1~Mini, with the electronics mounted on the back of the hand
(Figure~\ref{fig:enclosure}).

\begin{figure}[htbp]
\centering
\includegraphics[width=0.5\linewidth]{figures/mvp_hand.png}
\caption{The assembled four-finger Tactile Hand on its base, with a force-sensing
resistor bonded to each fingertip and its leads routed back toward the palm.
This four-finger morphology is the MVP platform used for all force-aware grasping
experiments.}
\label{fig:mvphand}
\end{figure}

\subsection{Sensing Electronics}
Each fingertip carries a Pololu \#2728 FSR. Because an FSR behaves as a
force-dependent variable resistor, we read it with a simple voltage divider: the
FSR is placed in series with a fixed reference resistor $R_\mathrm{ref}=22\,
\mathrm{k}\Omega$, and the ESP32-S3's 12-bit ADC samples the divider's midpoint.
The midpoint voltage is
\begin{equation}
V_\mathrm{adc} \;=\; V_\mathrm{cc}\,
\frac{R_\mathrm{ref}}{R_\mathrm{FSR}(F) + R_\mathrm{ref}},
\label{eq:divider}
\end{equation}
so as applied force $F$ increases, $R_\mathrm{FSR}$ falls and $V_\mathrm{adc}$
rises toward $V_\mathrm{cc}$. The ADC reports a raw count
$r \in [0, 4095]$. We chose the ESP32-S3 (WROOM-1
module)~\cite{espressif_esp32s3} for its four usable ADC channels, native USB,
and ample GPIO for servo control. To keep the analog signal clean we placed
100~nF and 10~$\mu$F decoupling capacitors near the ADC inputs, kept the analog
traces short, and sized the divider so the sensor operates in its most linear
band. The board is a solderable breadboard hosting the ESP32-S3, the four
voltage dividers, decoupling capacitors, and vertical 0.1-inch (2.54~mm)
headers for the FSR Dupont leads. Power comes from a 2S LiPo (7.4~V nominal)
stepped down to 5~V by a buck converter, with a Schottky diode preventing USB
backfeed during bench debugging.

\begin{figure}[htbp]
\centering
\includegraphics[width=0.95\linewidth]{figures/enclosure_circuit.png}
\caption{Custom ESP32-S3 sensing board (left) with the four FSR voltage-divider
channels highlighted, and the 3D-printed enclosure with status LEDs mounted on
the back of the hand (right).}
\label{fig:enclosure}
\end{figure}

\subsection{Firmware and ROS~2 Data Pipeline}
The firmware is built on the Arduino-ESP32 core. It samples the four ADC
channels at 100~Hz, applies the per-channel calibration described in
Section~\ref{sec:calibration} to convert raw counts to Newtons, and emits a
timestamped CSV line (timestamp followed by four force values) over a USB CDC
serial link at 115200~baud. During early bring-up we streamed synthetic force
values to validate the end-to-end data path before the real sensors were
attached, which let firmware and host software proceed in parallel with hardware
assembly.

On the host, a ROS~2 (Humble, on Ubuntu) serial-bridge node parses the CSV
stream and publishes the calibrated forces on \texttt{/tactile/forces}. The same
node subscribes to \texttt{/servo/commands} and forwards target servo positions
to the ESP32 over the same serial link. Publishing forces and commands as ROS~2
topics is the key architectural decision that let five subsystems
(electronics/firmware, calibration, simulation, vision, VR) integrate against a
single, stable interface.

\subsection{Force Sensor Calibration}
\label{sec:calibration}
FSRs are strongly nonlinear and vary unit to unit, so raw ADC counts must be
mapped to physical force. We built a calibration rig that holds a sensor flat
and level, then stacked 3D-printed PLA reference weights on a loading plate so a
known mass rests on the active area (Figure~\ref{fig:calibrig}). For each load we
logged the ADC reading over serial. Sweeping a range of loads produced a set of
(raw count, force) pairs.

\begin{figure}[htbp]
\centering
\includegraphics[width=0.55\linewidth]{figures/calib_rig.png}
\caption{Calibration rig. A 3D-printed loading plate rests on the FSR's active
area while stacked PLA reference weights apply a known mass; sweeping the stack
yields the (raw count, force) loading points used to fit the calibration curve.}
\label{fig:calibrig}
\end{figure}

The FSR's count-to-force relationship is well captured by a power law,
\begin{equation}
F(r) \;=\; a\, r^{\,b},
\label{eq:powerlaw}
\end{equation}
which we fit by nonlinear least squares. Across seven loading points the fit
yielded $a = 1.07\times10^{-9}$ and $b = 3.16$, with a root-mean-square error of
17.1~g and a maximum error of 25.7~g (Figure~\ref{fig:calib}). Since
$1\,\mathrm{N}\approx102$~g-force, a 25.7~g worst case corresponds to roughly
0.25~N---comfortably within our 1~N accuracy target. The procedure is
documented and repeatable; extending it to characterize all four fingers
individually (each with its own $a,b$) is in progress.

\begin{figure}[htbp]
\centering
\includegraphics[width=0.9\linewidth]{figures/analog.png}
\caption{Force-sensor calibration. A power-law model fit to seven stacked-weight
loading points maps raw ADC counts to applied force. Worst-case error (25.7~g
$\approx$ 0.25~N) is well within the 1~N target.}
\label{fig:calib}
\end{figure}

\subsection{Force-Aware Grasp Control}
\label{sec:graspctl}
The grasp controller (Figure~\ref{fig:block}) runs on the host at the sensor
rate. A finger is driven toward a target curl by an open-loop closing
trajectory (during MVP demos) or by a teleoperation input. At every cycle the
controller reads that finger's calibrated force from \texttt{/tactile/forces}.
While the force is below the safe threshold $\tau$, the commanded position is
forwarded to the servo and the finger continues to close. The instant the force
crosses $\tau$, the controller freezes that finger's target (the ``restrict
motion'' branch) so it stops squeezing harder even though the nominal trajectory would
keep closing, and a red status LED (one of two indicator LEDs on the board)
lights to signal that the threshold has been reached. Release is triggered on
demand by commanding the fingers open. In
practice this lets the hand wrap a thin-walled paper cup, halt as soon as the
fingertip force reaches the cup's safe limit, hold it stably, and release it
cleanly without buckling the wall (Figure~\ref{fig:grasp}). This closed loop is
the heart of the MVP, and it is where the tactile channel is actually used rather
than merely measured.

\begin{figure}[htbp]
\centering
\includegraphics[width=0.55\linewidth]{figures/closed_hand.png}
\caption{Force-aware grasping. The hand grasps a thin-walled cup and halts
finger closure when the measured fingertip force exceeds the safe threshold,
holding the cup without crushing it.}
\label{fig:grasp}
\end{figure}

\subsection{Vision-Based Hand-Pose Estimation}
To drive the hand from a human gesture we built a webcam pipeline on MediaPipe
Hands~\cite{zhang2020mediapipe} and OpenCV, which returns 21 hand landmarks in real time.
From the landmarks we estimate each finger's joint angles, reduce them to a
single normalized curl value in $[0,1]$ per finger, and map that value onto the
servo's angular range. The result is a four-element command vector published on
\texttt{/servo/commands}, so the robotic hand mirrors the operator's finger curl
live (Figure~\ref{fig:handpose}). Because the pipeline outputs the same
normalized curl interface used elsewhere, it is a drop-in alternative to the VR
input. Remaining work is robustness to lighting and partial occlusion and
recovering thumb abduction, which a single curl scalar does not capture.

\begin{figure}[htbp]
\centering
\includegraphics[width=0.95\linewidth]{figures/hand_pose.png}
\caption{Hand-pose mirroring. MediaPipe-based finger-pose estimation from a
webcam (right) is mapped to per-finger servo commands so the robotic hand
reproduces the operator's gesture (left).}
\label{fig:handpose}
\end{figure}

\subsection{Physics Simulation}
\label{sec:simsec}
We brought up the AmazingHand in the MuJoCo physics
simulator~\cite{todorov2012mujoco} by loading its
URDF/MJCF model and verifying that the simulated finger kinematics and curl
poses match the real hand (Figure~\ref{fig:sim}). The simulation serves two
purposes: it is a safe sandbox for validating control logic before running it on
hardware, and it is the foundation for sim-to-real experiments. The main one is an
in-hand cube-rotation policy trained in simulation that we intend to
transfer to the real hand, where the reality gap~\cite{zhao2020sim2real}
between simulated and physical FSR contact remains the central challenge.
Figure~\ref{fig:siminhand} shows the simulated hand mid-grasp on a cube, the
setup in which we are developing the in-hand cube-rotation policy, where
the simulated fingers reorient the object in place before we attempt to transfer
the learned behavior to the physical hand.

\begin{figure}[htbp]
\centering
\includegraphics[width=0.6\linewidth]{figures/sim_inhand.png}
\caption{Simulated in-hand manipulation. The AmazingHand model grasps a cube in
the physics simulator; this is the environment used to develop the in-hand
cube-rotation policy that we intend to transfer to hardware.}
\label{fig:siminhand}
\end{figure}

\begin{figure}[htbp]
\centering
\includegraphics[width=0.55\linewidth]{figures/sim.png}
\caption{The AmazingHand model loaded into MuJoCo, used to verify
finger kinematics and as groundwork for sim-to-real in-hand manipulation.}
\label{fig:sim}
\end{figure}

\subsection{VR Teleoperation and Arm Integration}
Our final input modality is VR teleoperation. We initially built a Meta Quest~2
APK in Unity 2022.3 LTS from the
\texttt{unity\_ros\_teleoperation} framework~\cite{leggedrobotics_teleop} and
successfully streamed 6-DoF controller pose (position and orientation) into
ROS~2 on Ubuntu. While functional, the Quest~2 path was awkward to integrate
with our Python-based control stack. We therefore switched to Valve
Index controllers, read through OpenVR/\texttt{pyopenvr}~\cite{valve2024openvr},
which gave us a clean pure-Python interface and native per-finger capacitive
tracking. With the Index controllers, controller pose is published as a ROS~2
topic at interactive rates. The same per-finger capacitive sensing lets the
operator's finger motion drive the hand's fingers directly
(Figure~\ref{fig:vrindex}), so a single controller supplies both the 6-DoF wrist
pose and a per-finger curl signal on the shared \texttt{/servo/commands}
interface.

\begin{figure}[htbp]
\centering
\includegraphics[width=0.85\linewidth]{figures/vr_index_hand.png}
\caption{Per-finger VR teleoperation. The Valve Index controller's native
capacitive finger tracking (right) is streamed over ROS~2 and mapped onto the
robotic hand's fingers (left), driving finger curl directly from the operator's
hand.}
\label{fig:vrindex}
\end{figure}

To turn streamed pose into motion we mounted the hand on an open-source ARCTOS
robotic arm~\cite{arctos2024} and mapped controller pose to the arm's
end-effector. Figure~\ref{fig:armmounted} shows the integrated platform: the
Tactile Hand is fixed to the ARCTOS end-effector, and the whole arm-hand chain
is driven from the VR controller through the same ROS~2 stack used for the
standalone hand. We built a full end-to-end teleoperation framework for
the arm: moving an Index controller moves the ARCTOS arm's end-effector through
our ROS~2 pipeline, with controller-to-arm frame calibration, inverse
kinematics, and motion smoothing all in place, so the full chain from VR input
to physical arm motion is closed (Figure~\ref{fig:vrarm}). The framework is not
yet production quality: we continue refining the frame calibration, jitter
filtering, latency, workspace-limit handling, and IK robustness toward reliable
teleoperated pick-and-place. Tightening those, and binding the per-finger Index
tracking to the force-aware hand controller, is the main remaining integration
task.

\begin{figure}[htbp]
\centering
\includegraphics[width=0.95\linewidth]{figures/arm_mounted.png}
\caption{Integrated hardware. The Tactile Hand mounted on the end-effector of
the open-source ARCTOS arm. The full arm-hand chain is driven from a VR
controller through the same ROS~2 stack used for the standalone hand.}
\label{fig:armmounted}
\end{figure}

\begin{figure}[htbp]
\centering
\includegraphics[width=0.58\linewidth]{figures/vr-demo.png}
\caption{VR teleoperation. Valve Index controller 6-DoF pose is streamed over
ROS~2 and mapped to the ARCTOS arm end-effector through an end-to-end framework
with frame calibration, IK, and motion smoothing, giving VR-driven control of
the arm-mounted hand that we continue to refine.}
\label{fig:vrarm}
\end{figure}

\subsection{Teleoperated Force-Aware Grasping}
\label{sec:teleopgrasp}
Mounting the hand on the arm let us ask a sharper question than the scripted MVP
demo could: does the force-aware safety guarantee survive once a human is in the
loop and the grasp is no longer a fixed trajectory? To test this we kept the
force-aware controller of Section~\ref{sec:graspctl} active on the arm-mounted
hand while an operator teleoperated the arm to a fixed, thin-walled 3D-printed
cup and triggered a grasp. We ran the same teleoperated grasp under two
conditions: once with the force-aware controller enabled, and once with force
feedback disabled so the fingers simply close to a fixed curl.
Figure~\ref{fig:teleopdemo} shows the full teleoperated setup under both
conditions: in each run the operator drives the same Index controller to the
same cup, and only the state of the force loop differs.

\begin{figure}[htbp]
\centering
\includegraphics[width=0.48\linewidth]{figures/teleop_noforce.png}\hfill
\includegraphics[width=0.48\linewidth]{figures/teleop_force.png}
\caption{Teleoperated grasp under both conditions. Left: with force feedback
disabled, the closing command drives the fingers into the cup with no stopping
criterion. Right: with force-threshold sensing enabled, each finger halts at the
cup's safe limit. The operator motion and the target object are the same in both
runs.}
\label{fig:teleopdemo}
\end{figure}

The contrast is stark and repeatable. With force-thresholding enabled, each
finger halts at the cup's safe limit and the hand holds the cup with its wall
intact; with feedback disabled, the identical teleoperated motion drives
the fingers past the safe force and visibly buckles the cup wall
(Figure~\ref{fig:teleopcompare}). Because nothing changes between the two runs
except the force loop, the comparison isolates the contribution of the tactile
channel: the outcome flips from a clean hold to a crushed cup. The same safety
behavior that protected the cup in the scripted MVP holds up when a human is
flying the arm in VR, so the loop is not confined to the scripted demo.

\begin{figure}[htbp]
\centering
\includegraphics[width=0.95\linewidth]{figures/teleop_compare.png}
\caption{Teleoperated force-aware grasping. Under the same VR-teleoperated
motion, the force-blind grasp (left) crushes the thin-walled cup, while the
force-threshold grasp (right) holds it with the wall intact. Only the force loop
differs between the two runs.}
\label{fig:teleopcompare}
\end{figure}

\subsection{Multi-Object Force-Awareness and Vision-Guided Grip Selection}
\label{sec:multiobject}
A single hand-tuned threshold is enough for one known object, but a genuinely
useful hand must choose an appropriate grip force for whatever it is handed. We
extended the force-aware controller toward automatic, per-object grip selection
and exercised it on a small set of everyday objects that demand different
grips. The hardest case is an empty versus a filled plastic bottle, which look
nearly identical but differ sharply in weight and therefore in the grip force
needed to hold them without slipping or crushing.

Our pipeline recognizes the target in a webcam frame with a fine-tuned
YOLO26 object detector~\cite{ultralytics_yolo26}. We annotated a small custom
dataset of our objects (e.g., filled plastic bottle, empty plastic
bottle, plastic box container, and book) and fine-tuned YOLO26 on
it, giving real-time bounding-box detections with a class label per object
(Figure~\ref{fig:yolo}). Each class maps to a hand-tuned target grip
force, with heavier or sturdier objects getting a firmer threshold and lighter or
more fragile ones a gentler threshold, and the detected class selects which
threshold the force-aware loop enforces. The detector changes only the
source of the target force; the downstream \texttt{/tactile/forces}
feedback path and the restrict-motion logic of Section~\ref{sec:graspctl} are
reused unchanged, which kept the addition lightweight. Growing the object set
and replacing the hand-tuned per-class thresholds with learned, weight-aware
force targets is left to future work.

\begin{figure}[htbp]
\centering
\includegraphics[width=0.8\linewidth]{figures/yolo_detection.png}
\caption{Vision-guided grip selection. A YOLO26 detector fine-tuned on a small
custom-annotated dataset recognizes each object and its class (e.g., filled vs.
empty plastic bottle); the detected class selects a per-object grip-force
threshold that the force-aware loop then enforces.}
\label{fig:yolo}
\end{figure}

This experiment also exposes a limitation of our calibration and motivates a
better sensor model. FSRs respond differently to off-axis and shear loading than
to the purely normal loading used during calibration
(Section~\ref{sec:calibration}), so a grip that is slightly misaligned reads a
biased force. We have begun a system-identification effort that characterizes
each sensor's response to off-axis force and feeds the measured response back
into the physics simulation (Section~\ref{sec:simsec}), with the goal of
tightening the sim-to-real match for the contact-rich grasps these multi-object
experiments demand.

%=====================================================================
\section{Milestones}
\label{sec:milestones}
%=====================================================================
We targeted MVP completion by Week~6, with reach goals in Weeks~7--10.
Table~\ref{tab:milestones} lists each milestone and its final status: all
MVP-track items were completed on schedule, and the reach goals range from
complete (multi-object grip selection; teleoperated force-aware grasp) through
demonstrated-but-unrefined (VR-to-arm teleoperation) to not completed
(sim-to-real transfer).

\begin{table}[t]
\caption{Milestone status. All MVP-track items are complete; reach goals range
from complete to not completed, consistent with their deferred priority.}
\label{tab:milestones}
\small
\renewcommand{\arraystretch}{1.2}
\begin{tabularx}{\linewidth}{@{}Xll@{}}
\toprule
\textbf{Milestone} & \textbf{Target} & \textbf{Status} \\
\midrule
System design \& component selection      & W2--W3 & Complete \\
Prototype board design \& fabrication     & W3--W4 & Complete \\
Hand assembly \& sensor integration       & W3--W5 & Complete \\
ESP32 firmware: read FSRs, stream USB     & W4--W5 & Complete \\
FSR calibration to within 1~N             & W4--W6 & Complete \\
Simulation bring-up of AmazingHand        & W3--W5 & Complete \\
Force-aware grasping closed loop          & W5--W6 & Complete \\
\textbf{MVP completion \& validation}     & \textbf{W6} & \textbf{Complete} \\
Hand-pose estimation pipeline (reach)     & W7--W9 & Complete \\
VR controller pose streaming (reach)      & W7--W9 & Complete \\
Teleoperated force-aware grasp demo (reach)& W9    & Complete \\
VR $\rightarrow$ arm end-effector (reach) & W8--W9 & Demonstrated \\
Multi-object force awareness (reach)      & W8--W9 & Complete \\
Per-finger calibration table (reach)      & W8     & In progress \\
Sim-to-real in-hand rotation (reach)      & W8--W9 & Not completed \\
Final paper, demo video, GitHub repo      & W9--W10& In progress \\
\bottomrule
\end{tabularx}
\end{table}

\subsection{Mid-Quarter Revisions}
We revised the plan twice. After feedback on our specification, we tied every
subsystem to a qualified owner and rewrote the milestone table and Gantt chart to
separate gating items from parallel work, making schedule risk easy to read. We
also changed the VR hardware: the charter specified a Meta Quest~2, and we did
deploy a working Quest~2 APK, but integrating it with our Python stack was
cumbersome, so we switched to Valve Index controllers
(OpenVR/\texttt{pyopenvr}) for a cleaner pure-Python interface and native
per-finger tracking.

\subsection{What We Did and Did Not Finish}
All MVP deliverables were met. Among reach goals, multi-object grip selection was
completed (YOLO26 detection driving per-class force thresholds,
Section~\ref{sec:multiobject}) and the teleoperated force-aware grasp was
demonstrated (Section~\ref{sec:teleopgrasp}). VR-to-arm teleoperation works end
to end through a full framework with frame calibration, inverse kinematics, and
motion smoothing, which we continue to refine (jitter filtering, latency, and IK
robustness) toward reliable pick-and-place. Three items remain unfinished: the sim-to-real transfer of the
in-hand rotation policy (trained in simulation but not yet moved to hardware),
finger adduction/abduction control (our curl interface captures flexion only),
and per-finger calibration across all four sensors (validated on one sensor,
being repeated per finger). Because reach goals were started only after the
Week~6 MVP, these gaps reflect the time available afterward, not a failure of the
core deliverable.

\subsection{Problems and Mitigations}
FSR nonlinearity and unit-to-unit variation threatened our 1~N target; the
stacked-weight calibration rig and power-law model brought worst-case error to
0.25~N. Analog noise was handled with decoupling capacitors and short traces, and
the Quest~2 integration friction by switching to Valve Index. Mounting the hand
on the arm surfaced three further issues: ARCTOS wiring needed rework before the
arm drove reliably; off-axis fingertip loading during arm-driven grasps biases
readings relative to the normal-load calibration (we are characterizing this and
feeding it back into simulation, Section~\ref{sec:multiobject}); and the
read--decide--actuate loop has enough latency that the operator must close slowly
to avoid overshooting the threshold. The hand's wide palm limited its
reachability and made small objects hard to grasp, so we favored larger test
objects. Throughout, the MVP-first policy---integrate reach goals
only after the MVP, and replan any task slipping more than a week---kept the
gating path protected.

%=====================================================================
\section{Conclusion}
\label{sec:conclusion}
%=====================================================================
A robot hand that cannot feel how hard it is gripping is dangerous, and until
now that safety feature has been missing from affordable open-source hands. We
presented Tactile Hands, a sub-\$200 four-fingered robotic hand that adds
calibrated fingertip force sensing to the open AmazingHand platform. The system
reads four FSRs through a custom ESP32-S3 board, converts raw readings to force
with better than 1N accuracy using a power-law calibration, and runs a
force-aware controller that grasps and releases a fragile object without
crushing it. We additionally demonstrated webcam-based hand-pose mirroring,
physics simulation of the hand, and VR controller streaming over ROS~2, we
mounted the hand on a VR-teleoperated arm, and we added a fine-tuned YOLO26
detector that selects a per-object grip force automatically. On that arm-mounted
configuration we showed our central result: under identical teleoperated motion,
the force-aware grasp holds a thin-walled cup while a force-blind grasp crushes
it. The safety guarantee carries over to teleoperation instead of being limited
to a scripted demo.

The significance of the work is accessibility. It shows that the most
safety-critical sensing capability of a robot hand, knowing its own grip
force, can be added to a cheap, reproducible platform without specialized
hardware, and that it keeps working once the hand is placed on an arm and driven
by a human, lowering the barrier to safe manipulation research and education.

Future work follows directly from our partial milestones. We will complete
per-finger calibration so each fingertip carries its own coefficients; expand
the vision-guided, multi-object grip-selection pipeline by growing the YOLO26
object dataset and replacing the hand-tuned per-class thresholds with learned,
weight-aware force targets; carry the
off-axis system-identification results into the simulator to shrink the
sim-to-real gap and push the in-hand cube-rotation policy onto real hardware;
and, most substantially, refine our end-to-end VR arm teleoperation into
reliable pick-and-place by tightening the controller-to-arm calibration,
filtering controller motion, and binding the Index per-finger tracking to the
force-aware hand controller. All hardware designs, firmware, calibration data, and software
will be released in a public GitHub repository so that others can reproduce a
low-cost, force-sensing dexterous hand.

%=====================================================================
\bibliographystyle{ACM-Reference-Format}
\bibliography{references}

\end{document}